%===========================================================
% Capítulo: Desarrollo — Ingeniería del pipeline, organización y reproducibilidad
%===========================================================
\chapter{Desarrollo}
\label{chap:desarrollo}

\section*{Visión general del pipeline}
El sistema se implementa como una cadena determinista y modular que va desde los hologramas crudos hasta una decisión a nivel de sujeto. La cadena consta de: (i) ingestión y control de calidad (QC) con trazabilidad completa; (ii) preprocesamiento ligero físicamente motivado (percentiles, CLAHE y estandarización) para estabilizar la estructura interferométrica; (iii) muestreo reproducible de parches, con filtros por energía y SNR espectral; (iv) extracción de rasgos físico–estadísticos (LBP, GLCM, FFT y anisotropía) definidos en fracción de Nyquist para mitigar dependencias instrumentales; (v) agregación parche $\rightarrow$ sujeto mediante funcionales robustas (media, logit–media, softmax-$\beta$); (vi) entrenamiento y evaluación con \emph{validación cruzada anidada} estratificada por sujeto; (vii) calibración probabilística y selección de umbral por costo clínico; y (viii) interpretabilidad (\emph{post-hoc}) para alinear las importancias con la física de difracción. El diseño prioriza parquedad, interpretabilidad y reproducibilidad, con costes de cómputo compatibles con despliegue en CPU de bajo consumo \cite{Chen2022Photonics,Greenbaum2014STM,Tsamardinos2022Patterns}.

\section*{Arquitectura, trazabilidad y control de versiones}
La organización del proyecto separa datos crudos, artefactos intermedios, rasgos procesados, configuraciones y registros. Los hologramas en \texttt{data/raw} son inmutables; toda transformación produce artefactos versionados en \texttt{data/interim} (parches y métricas de QC) y \texttt{data/processed} (rasgos por parche y por sujeto). Cada corrida registra semillas, versiones de librerías, rutas, hashes de integridad, particiones \emph{a priori} y parámetros de búsqueda. La configuración es declarativa (archivos YAML) y controla tanto hiperparámetros de preprocesamiento y rasgos como del clasificador y la validación. Esta separación minimiza el ``\emph{configuration drift}'' y facilita ablasiones sistemáticas.

\section*{Particiones honestas y prevención de fugas}
La unidad estadística es el \textbf{sujeto/slide}. Toda partición se hace con \emph{GroupKFold} (grupos = \texttt{subject\_id}), de forma que ninguna vista del mismo sujeto aparece simultáneamente en entrenamiento y prueba externa. Sobre cada pliegue externo, la selección de hiperparámetros se realiza en un bucle interno independiente, maximizando AUC a nivel de \emph{sujeto} tras la agregación parche $\rightarrow$ sujeto. De este modo se separa estrictamente la selección del modelo de su evaluación, evitando el optimismo de validaciones no anidadas \cite{Tsamardinos2022Patterns}. Las predicciones \emph{out-of-fold} concatenadas constituyen la base para estimar AUC/AUPRC e Intervalos de Confianza (IC95\%) por \emph{bootstrap de sujetos}.

\section*{Control de calidad (QC) y homogeneización de intensidad}
El QC descarta campos con saturación global, desenfoque severo (baja varianza del Laplaciano), artefactos extensos o exposiciones anómalas. Todos los hologramas se convierten a escala de grises \texttt{float32} y se normalizan por el rango dinámico nativo del sensor, registrando metadatos relevantes (tamaño de píxel $p$, longitud de onda $\lambda$ y distancia $z$ si está disponible). Esta homogeneización prepara la señal para un preprocesamiento que respeta la estructura interferométrica y las restricciones de muestreo discutidas en el marco teórico \cite{Chen2022Photonics}.

\section*{Preprocesamiento determinista y físicamente motivado}
La normalización por percentiles (1–99) remueve colas extremas y estabiliza el rango dinámico sin sacrificar las franjas centrales del patrón holográfico. CLAHE, con \emph{clip-limit} y tamaño de tesela seleccionados por validación interna, amplifica el microcontraste local donde las franjas tienen baja visibilidad, pero limita la saturación, preservando la estructura periódica de interferencia \cite{Zhou2019MedImageContrast}. Finalmente, la estandarización $z$-score por imagen centra y escala intensidades para que descriptores como LBP/GLCM respondan a diferencias de \emph{patrón} y no a \emph{offsets} o ganancias globales. Esta etapa es determinista y encapsulada, garantizando que estadísticas de escalado y cuantización se estiman exclusivamente sobre el entrenamiento del pliegue correspondiente.

\section*{Muestreo de parches y filtros energía/SNR}
Para evitar sesgos de detectores de célula y mantener reproducibilidad, se adopta una rejilla uniforme de parches de tamaño $w\times w$ con \emph{stride} fijo, acotado para cubrir varias franjas alrededor del máximo central típico. El tamaño $w$ se elige para capturar suficiente estructura (al menos 2–3 anillos del primer lóbulo), con sensibilidad de hiperparámetro explorada en la búsqueda interna. Se aplican dos filtros simples y efectivos: (i) un umbral de energía espectral total (percentil bajo del entrenamiento) para eliminar parches triviales, y (ii) una SNR espectral local basada en el perfil radial en una banda anular $B$ que captura el primer lóbulo. Estos filtros reducen variabilidad no informativa y estabilizan la distribución de rasgos sin introducir heurísticas dependientes de segmentación.

\section*{Extracción de rasgos robustos e invariantes}
El vector de rasgos concatena tres familias complementarias, cada una con fundamento físico y estadístico. 
\emph{(1) Textura LBP multiescala:} códigos uniformes en radios $R$ y números de vecinos $P$ combinados con momentos del histograma (energía, entropía, contraste) caracterizan la granularidad periódica de las franjas; en AF se espera mayor entropía y contraste por el enriquecimiento de altas frecuencias y la presencia de anisotropías finas. 
\emph{(2) Co-ocurrencia GLCM:} energía, contraste, homogeneidad, correlación y entropía promediados sobre distancias y ángulos cuantifican orden/desorden interferométrico de segundo orden; en AF suelen aumentar el contraste y la entropía y disminuir la homogeneidad \cite{Kiani2022TextureSurvey}. 
\emph{(3) Espectro de Fourier y anisotropía:} el perfil radial en fracción de Nyquist, las posiciones de picos $(k_1,k_2)$, la energía por bandas y la razón $R_{HL}=E_\mathrm{hi}/E_\mathrm{lo}$ capturan la escala de franjas (ligada a tamaño/elongación efectiva del objeto holografiado). Medidas angulares (\emph{varianza} y modos periódicos $m=2,4$) sintetizan el grado de elipticidad/direccionalidad del patrón. Definir todas las frecuencias con respecto a $k_\mathrm{NQ}=\pi/p$ mitiga el \emph{domain shift} entre sensores \cite{Chen2022Photonics}.

\section*{Agregación parche $\rightarrow$ sujeto y nivel clínico de decisión}
La inferencia clínica se decide a nivel de sujeto o slide. Dado un conjunto de puntajes por parche, se emplean agregadores robustos: la media (estable ante ruido), la logit–media (que promedia en el espacio de \emph{log-odds} y reduce saturaciones), o un softmax-$\beta$ (que aproxima el máximo y es útil cuando pocos parches positivos bastan para indicar AF). La selección del agregador se incluye como hiperparámetro en la validación interna, y todas las métricas de evaluación externa se computan \emph{a nivel de sujeto} para evitar la inflación de desempeño por múltiples parches correlacionados.

\section*{Selección de rasgos, parquedad e interpretabilidad}
Para favorecer parquedad y estabilidad, se utiliza un filtro univariante con control de FDR (Benjamini–Hochberg) que retiene un conjunto de rasgos candidatos, seguido de \emph{importancias por permutación} y/o \emph{RFE} específicas del clasificador. La hipótesis H4 anticipa que un subconjunto compacto ($\leq 10$ rasgos) preserva la AUC con caída menor a 0.02. Esta parquedad facilita auditoría y reduce varianza, y sus importancias se contrastan con expectativas físicas: picos radiales más altos, mayor razón $R_{HL}$ y mayor entropía/contraste en LBP/GLCM son coherentes con patrones más finos y anisotrópicos de AF.

\section*{Modelado y búsqueda de hiperparámetros}
Se comparan modelos parcos e interpretables (Regresión Logística $\ell_2$, SVM-RBF con probabilidades calibradas y Random Forest con importancias por permutación) y, como línea base, una CNN ligera con transferencia para establecer un techo de referencia en el régimen de pocos datos \cite{Alzubaidi2020Electronics}. La búsqueda de hiperparámetros se ejecuta en el bucle interno de la validación anidada con rejillas logarítmicas razonables (p.\,ej., $C$ y $\gamma$ en SVM, profundidad y número de árboles en RF), siempre evaluando AUC \emph{a nivel de sujeto} tras la agregación.

\section*{Validación cruzada anidada e intervalos de confianza}
Cada pliegue externo produce predicciones \emph{out-of-fold} a nivel de sujeto, que se concatenan para estimar AUC/AUPRC imparciales. La incertidumbre se cuantifica mediante \emph{bootstrap} de sujetos (miles de réplicas) reportando IC95\% (percentil o BCa). Este enfoque captura la variabilidad inter–sujeto (relevante clínicamente) y evita el sesgo de selección, cumpliendo recomendaciones metodológicas para tamaños muestrales modestos \cite{Tsamardinos2022Patterns}.

\section*{Calibración probabilística y umbral clínico}
Las probabilidades calibradas se verifican con \emph{Brier score} y curvas de confiabilidad. Si procede, se aplica \emph{Platt scaling} o \emph{regresión isotónica} usando exclusivamente las predicciones \emph{out-of-fold} del bucle externo (evitando doble uso de datos). La decisión binaria se fija por el umbral que minimiza el riesgo esperado bajo una matriz de costos clínicos; típicamente $C_\mathrm{FN}\gg C_\mathrm{FP}$ en cribado, lo que conduce a umbrales bajos (alta sensibilidad). Alternativamente, se reporta el punto de Youden como compromiso estándar. Se presentan matrices de confusión, sensibilidad, especificidad, PPV/NPV y F1 en ambos puntos operativos.

\section*{Interpretabilidad y validación física de las explicaciones}
Las explicaciones \emph{post-hoc} (TreeSHAP para RF y KernelSHAP para SVM/LR) descomponen las predicciones en contribuciones por rasgo. Se espera que los rasgos que más empujan hacia la clase AF sean físicamente plausibles: incremento del primer pico radial $k_1$, mayor razón $R_{HL}$, pendiente espectral menos negativa y aumento de entropía/contraste en LBP/GLCM. Se incluyen resúmenes globales (\emph{beeswarm}) y análisis de casos límite (\emph{waterfall}) por pliegue externo, verificando coherencia con las hipótesis H1–H2 y el modelo de difracción \cite{Kiani2022TextureSurvey,Chen2022Photonics,Greenbaum2014STM}.

\section*{Ablaciones y análisis de sensibilidad}
Se realizan ablasiones controladas manteniendo la misma validación anidada y calibración: LBP vs.\ GLCM vs.\ FFT, combinaciones por pares, y la triada completa; también se explora la sensibilidad a \emph{clip-limit} y tamaño de tesela (CLAHE), niveles de cuantización para GLCM, tamaño/stride de parche y elección de agregador. La robustez a SNR se evalúa estratificando por cuartiles de $\mathrm{SNR}_B$ y comparando AUC/IC por estrato; una caída $< 5$ puntos entre Q1 y Q4 apoya la robustez del diseño.

\section*{Complejidad computacional e inferencia en CPU}
Con parches de $w\in[96,128]$, el coste dominante por parche es la FFT 2D ($O(w^2\log w)$), seguida por LBP/GLCM con parámetros acotados; la extracción completa y la inferencia del clasificador se mantienen por debajo del milisegundo por parche en CPU moderna. Con del orden de $10^2$ parches por sujeto, la decisión a nivel de slide se obtiene en $\leq 0.1$\,s. El almacenamiento de rasgos es modesto ($\mathcal{O}(10^2)$ valores en \texttt{float32} por parche). Estos presupuestos satisfacen el criterio de eficiencia (C4) y son compatibles con despliegues en equipos portátiles o \emph{edge devices}.

\section*{Buenas prácticas de reproducibilidad}
Todas las transformaciones se encapsulan en \emph{pipelines} inmutables con semillas fijadas; se guardan configuraciones completas de cada corrida, predicciones \emph{out-of-fold}, curvas ROC/PR y calibradores ajustados por pliegue. El reporte de métricas incluye IC95\% y, cuando es pertinente, bandas de confianza para curvas. La documentación detalla versiones de librerías y hardware utilizado, permitiendo auditoría y réplica.

\section*{Gestión de riesgos y mitigaciones}
Se anticipan tres riesgos principales: (i) \emph{sobreajuste} por pocos sujetos, mitigado por validación anidada, parquedad de rasgos y reporte con IC; (ii) \emph{artefactos/gemela} que distorsionen patrones, mitigados por filtros energía/SNR y el uso de descriptores en el dominio de Fourier y co-ocurrencia que sintetizan el efecto sin reconstrucción; y (iii) \emph{desplazamientos de dominio} por cambios en $p$, $\lambda$ o $z$, atenuados al expresar frecuencias en fracción de Nyquist, adoptar rasgos multiescala y realizar análisis de sensibilidad. La comparación con una CNN ligera se mantiene como control metodológico más que como solución principal interpretativa \cite{Alzubaidi2020Electronics}.

\section*{Despliegue y consideraciones prácticas (POCT)}
Para el uso en punto de atención, el pipeline puede empaquetarse en una aplicación de escritorio ligera o un contenedor con dependencias mínimas, procesando en lote campos completos y emitiendo una \emph{probabilidad calibrada} por sujeto junto con explicaciones de rasgos dominantes. La compatibilidad con plataformas LDHM de bajo costo ya demostradas en la literatura facilita la transferencia a hardware económico, siempre que se respeten las condiciones de muestreo y coherencia discutidas \cite{Amann2019SciRep,BuitragoDuque2023HardwareX,Greenbaum2014STM}. Se recomiendan chequeos automáticos de QC y una bitácora de calibración periódica.

\section*{Resultados esperados y criterios de aceptación}
El sistema se considera listo para validación externa si, en la validación anidada interna, cumple: AUC-ROC media $\ge 0.90$ con IC95\% por encima de 0.80; sensibilidad $\ge 0.90$ con especificidad $\ge 0.85$ en el umbral clínico; estabilidad inter–pliegue (DE de AUC $\le 0.05$) y parquedad (subconjunto $\le 10$ rasgos con $\Delta\text{AUC}\le 0.02$). Estos criterios enlazan directamente con los objetivos y las hipótesis planteadas, y proporcionan una base transparente para auditoría y toma de decisiones clínicas.

\section*{Síntesis}
El desarrollo implementa fielmente el marco físico y estadístico: los hologramas, tratados como señales interferométricas, se proyectan a un espacio de rasgos que captura escala, direccionalidad y orden local de franjas; los modelos parcos y la validación anidada garantizan estimaciones honestas en regímenes de pocos datos; y la calibración más interpretabilidad cierran el ciclo hacia decisiones clínicas informadas. La arquitectura, reproducible y eficiente, es compatible con prototipos LDHM de bajo costo, apoyando el objetivo de un tamizaje portable para anemia falciforme \cite{Amann2019SciRep,BuitragoDuque2023HardwareX,Greenbaum2014STM,Chen2022Photonics,Kiani2022TextureSurvey,Zhou2019MedImageContrast,Tsamardinos2022Patterns,Alzubaidi2020Electronics}.
